% Options for packages loaded elsewhere
% Options for packages loaded elsewhere
\PassOptionsToPackage{unicode}{hyperref}
\PassOptionsToPackage{hyphens}{url}
\PassOptionsToPackage{dvipsnames,svgnames,x11names}{xcolor}
%
\documentclass[
  letterpaper,
]{article}
\usepackage{xcolor}
\usepackage{amsmath,amssymb}
\setcounter{secnumdepth}{5}
\usepackage{iftex}
\ifPDFTeX
  \usepackage[T1]{fontenc}
  \usepackage[utf8]{inputenc}
  \usepackage{textcomp} % provide euro and other symbols
\else % if luatex or xetex
  \usepackage{unicode-math} % this also loads fontspec
  \defaultfontfeatures{Scale=MatchLowercase}
  \defaultfontfeatures[\rmfamily]{Ligatures=TeX,Scale=1}
\fi
\usepackage{lmodern}
\ifPDFTeX\else
  % xetex/luatex font selection
\fi
% Use upquote if available, for straight quotes in verbatim environments
\IfFileExists{upquote.sty}{\usepackage{upquote}}{}
\IfFileExists{microtype.sty}{% use microtype if available
  \usepackage[]{microtype}
  \UseMicrotypeSet[protrusion]{basicmath} % disable protrusion for tt fonts
}{}
\makeatletter
\@ifundefined{KOMAClassName}{% if non-KOMA class
  \IfFileExists{parskip.sty}{%
    \usepackage{parskip}
  }{% else
    \setlength{\parindent}{0pt}
    \setlength{\parskip}{6pt plus 2pt minus 1pt}}
}{% if KOMA class
  \KOMAoptions{parskip=half}}
\makeatother
% Make \paragraph and \subparagraph free-standing
\makeatletter
\ifx\paragraph\undefined\else
  \let\oldparagraph\paragraph
  \renewcommand{\paragraph}{
    \@ifstar
      \xxxParagraphStar
      \xxxParagraphNoStar
  }
  \newcommand{\xxxParagraphStar}[1]{\oldparagraph*{#1}\mbox{}}
  \newcommand{\xxxParagraphNoStar}[1]{\oldparagraph{#1}\mbox{}}
\fi
\ifx\subparagraph\undefined\else
  \let\oldsubparagraph\subparagraph
  \renewcommand{\subparagraph}{
    \@ifstar
      \xxxSubParagraphStar
      \xxxSubParagraphNoStar
  }
  \newcommand{\xxxSubParagraphStar}[1]{\oldsubparagraph*{#1}\mbox{}}
  \newcommand{\xxxSubParagraphNoStar}[1]{\oldsubparagraph{#1}\mbox{}}
\fi
\makeatother


\usepackage{longtable,booktabs,array}
\usepackage{calc} % for calculating minipage widths
% Correct order of tables after \paragraph or \subparagraph
\usepackage{etoolbox}
\makeatletter
\patchcmd\longtable{\par}{\if@noskipsec\mbox{}\fi\par}{}{}
\makeatother
% Allow footnotes in longtable head/foot
\IfFileExists{footnotehyper.sty}{\usepackage{footnotehyper}}{\usepackage{footnote}}
\makesavenoteenv{longtable}
\usepackage{graphicx}
\makeatletter
\newsavebox\pandoc@box
\newcommand*\pandocbounded[1]{% scales image to fit in text height/width
  \sbox\pandoc@box{#1}%
  \Gscale@div\@tempa{\textheight}{\dimexpr\ht\pandoc@box+\dp\pandoc@box\relax}%
  \Gscale@div\@tempb{\linewidth}{\wd\pandoc@box}%
  \ifdim\@tempb\p@<\@tempa\p@\let\@tempa\@tempb\fi% select the smaller of both
  \ifdim\@tempa\p@<\p@\scalebox{\@tempa}{\usebox\pandoc@box}%
  \else\usebox{\pandoc@box}%
  \fi%
}
% Set default figure placement to htbp
\def\fps@figure{htbp}
\makeatother


% definitions for citeproc citations
\NewDocumentCommand\citeproctext{}{}
\NewDocumentCommand\citeproc{mm}{%
  \begingroup\def\citeproctext{#2}\cite{#1}\endgroup}
\makeatletter
 % allow citations to break across lines
 \let\@cite@ofmt\@firstofone
 % avoid brackets around text for \cite:
 \def\@biblabel#1{}
 \def\@cite#1#2{{#1\if@tempswa , #2\fi}}
\makeatother
\newlength{\cslhangindent}
\setlength{\cslhangindent}{1.5em}
\newlength{\csllabelwidth}
\setlength{\csllabelwidth}{3em}
\newenvironment{CSLReferences}[2] % #1 hanging-indent, #2 entry-spacing
 {\begin{list}{}{%
  \setlength{\itemindent}{0pt}
  \setlength{\leftmargin}{0pt}
  \setlength{\parsep}{0pt}
  % turn on hanging indent if param 1 is 1
  \ifodd #1
   \setlength{\leftmargin}{\cslhangindent}
   \setlength{\itemindent}{-1\cslhangindent}
  \fi
  % set entry spacing
  \setlength{\itemsep}{#2\baselineskip}}}
 {\end{list}}
\usepackage{calc}
\newcommand{\CSLBlock}[1]{\hfill\break\parbox[t]{\linewidth}{\strut\ignorespaces#1\strut}}
\newcommand{\CSLLeftMargin}[1]{\parbox[t]{\csllabelwidth}{\strut#1\strut}}
\newcommand{\CSLRightInline}[1]{\parbox[t]{\linewidth - \csllabelwidth}{\strut#1\strut}}
\newcommand{\CSLIndent}[1]{\hspace{\cslhangindent}#1}



\setlength{\emergencystretch}{3em} % prevent overfull lines

\providecommand{\tightlist}{%
  \setlength{\itemsep}{0pt}\setlength{\parskip}{0pt}}



 


\usepackage{hyphenat}
\usepackage{graphicx}
% and their extensions so you won't have to specify these with
 % every instance of \includegraphics
\usepackage{pdfcomment}
\DeclareGraphicsExtensions{.pdf,.jpeg,.png}
\usepackage{wallpaper} % for the background image on title page
\usepackage{geometry}

\usepackage{lastpage}
% set font
% \usepackage{fontspec}
% \setsansfont{Latin Modern Sans}
% \renewcommand{\setmainfont}[2][]{\fontspec[#1]{Latin Modern Sans}}
% \renewcommand{\rmdefault}{lmss}

% Acronyms
\usepackage[acronym]{glossaries}
\glsdisablehyper
\makenoidxglossaries
\loadglsentries{report_glossary.tex}

% Font setup (LuaLaTeX/XeLaTeX via Quarto)
\usepackage{fontspec}
\defaultfontfeatures{Ligatures=TeX}

% Explicit face mapping for Latin Modern Sans
\setsansfont{Latin Modern Sans}[]

% Make entire document sans-serif
\setmainfont{Latin Modern Sans}[]
\renewcommand{\familydefault}{\sfdefault}

% \usepackage{lmodern}
% \usepackage[T1]{fontenc}

\newfontfamily\sectionfont[Color=Black]{Latin Modern Sans}
\newfontfamily\subsectionfont[Color=Black]{Latin Modern Sans}
\newfontfamily\subsubsectionfont[Color=Black]{Latin Modern Sans}
% \addtokomafont{section}{\sectionfont}
% \addtokomafont{subsection}{\subsectionfont}
% \addtokomafont{subsubsection}{\subsubsectionfont}
\usepackage[headsepline=0.005pt:,footsepline=0.005pt:,plainfootsepline,automark]{scrlayer-scrpage}
\clearpairofpagestyles
\ohead[]{\headmark} \cofoot[\pagemark]{\pagemark}
\lohead{Category 3 stocks assessment 2023}
\ModifyLayer[addvoffset=-.6ex]{scrheadings.foot.above.line}
\ModifyLayer[addvoffset=-.6ex]{plain.scrheadings.foot.above.line}
\setkomafont{pageheadfoot}{\small}

% add soul package to remove latex error
\usepackage{soul}

\usepackage{float}
\floatplacement{table}{H}

\makeatletter
\@ifpackageloaded{caption}{}{\usepackage{caption}}
\AtBeginDocument{%
\ifdefined\contentsname
  \renewcommand*\contentsname{Table of contents}
\else
  \newcommand\contentsname{Table of contents}
\fi
\ifdefined\listfigurename
  \renewcommand*\listfigurename{List of Figures}
\else
  \newcommand\listfigurename{List of Figures}
\fi
\ifdefined\listtablename
  \renewcommand*\listtablename{List of Tables}
\else
  \newcommand\listtablename{List of Tables}
\fi
\ifdefined\figurename
  \renewcommand*\figurename{Figure}
\else
  \newcommand\figurename{Figure}
\fi
\ifdefined\tablename
  \renewcommand*\tablename{Table}
\else
  \newcommand\tablename{Table}
\fi
}
\@ifpackageloaded{float}{}{\usepackage{float}}
\floatstyle{ruled}
\@ifundefined{c@chapter}{\newfloat{codelisting}{h}{lop}}{\newfloat{codelisting}{h}{lop}[chapter]}
\floatname{codelisting}{Listing}
\newcommand*\listoflistings{\listof{codelisting}{List of Listings}}
\makeatother
\makeatletter
\makeatother
\makeatletter
\@ifpackageloaded{caption}{}{\usepackage{caption}}
\@ifpackageloaded{subcaption}{}{\usepackage{subcaption}}
\makeatother
\usepackage{bookmark}
\IfFileExists{xurl.sty}{\usepackage{xurl}}{} % add URL line breaks if available
\urlstyle{same}
\hypersetup{
  pdftitle={Status and Catch Limit Estimation of Category 3 stocks off the U.S. West Coast in 2027},
  pdfauthor={Jason M. Cope; Chantel R. Wetzel},
  colorlinks=true,
  linkcolor={blue},
  filecolor={Maroon},
  citecolor={Blue},
  urlcolor={Blue},
  pdfcreator={LaTeX via pandoc}}


\title{Status and Catch Limit Estimation of Category 3 stocks off the
U.S. West Coast in 2027}
\author{Jason M. Cope \and Chantel R. Wetzel}
\date{2026-11-25}
\begin{document}
  \begin{titlepage}
  % This is a combination of Pandoc templating and LaTeX
  % Pandoc templating https://pandoc.org/MANUAL.html#templates
  % See the README for help

  \newgeometry{top=2in,bottom=1in,right=1in,left=1in}
  \begin{minipage}[b][\textheight][s]{\textwidth}
  % Ross would've subbed lines 6, 8 with these lines:
  %\newgeometry{top=2in,bottom=1in,right=1in,left=1in}%
  %\noindent  %\tracingall
  %\begin{minipage}[b][\textheight][s]{.975\textwidth}%% RRM: avoid Overfull box


  \raggedright

  % \includegraphics[width=2cm]{NOAA_Transparent_Logo.png}

  % background image


  % Title and subtitle
  {\huge\bfseries\nohyphens{Status and Catch Limit Estimation of
  Category 3 stocks off the U.S. West Coast in 2027}}\\[1\baselineskip]
  % Ross would change the end of the above line to the following because \par must come before the group closes and line-depth reverts.
  % }\par}%\\[1\baselineskip]



  \vspace{2\baselineskip}
  % Ross would change this to 2\baselineskip

  %%%%%% Cover image
  \pandocbounded{\includegraphics[width=6in, alt=An illustration of Category 3 stocks.]{support\_files/Category\_3\_stocks.png}}
  % cover page customization need to be inserted here
  %\pdftooltip{\includegraphics{support\_files/Category\_3\_stocks.png}}{Alt text}

  \vspace{1\baselineskip}

  % Authors
  % This hairy bit of code is just to get "and" between the last 2
  % authors. See below if you don't need that
   {\large{Jason M. Cope}}{\textsuperscript{1}}%
  %
  %
  { and \large{Chantel R. Wetzel}}%
  {\textsuperscript{1}}%
  %


  % This is how to do it if you don't need the "and"

  %%%%%% Affiliations
  \vspace{2\baselineskip}

  \hangindent=1em
  \hangafter=1
  % Ross would change the above line to:
  % \hangafter=1\relax
  %
  {1}.~{NOAA Fisheries Northwest Fisheries Science Center}%
  %
  %
    , %
    2725 Montlake Boulevard East, Seattle, WA, 98112-2097%
  %


  %%%%%% Correspondence
  \vspace{1\baselineskip}


  %use \vfill instead to get the space to fill flexibly
  %\vspace{0.25\textheight} % Whitespace between the title block and the publisher

  \vfill


  % Whitespace between the title block and the tagline
  \vspace{1\baselineskip}

  %%%%%% Tagline at bottom
  % Ross says the tagline below could also be centered
  \includegraphics[width=2cm, artifact]{support_files/us_doc_logo.png}\newline % empty curly brackets without alt text is suitable for this logo because it's purely decorative/an "artifact"
  U.S. Department of Commerce\newline
  National Oceanic and Atmospheric Administration\newline
  National Marine Fisheries Service\newline
  Northwest Fisheries Science Center\newline

  \end{minipage}
  \restoregeometry
  \end{titlepage}

\newgeometry{top=1in,bottom=1in,right=1in,left=1in}

\pdfbookmark[1]{Table of Contents}{toc}
\renewcommand*\contentsname{Table of contents}
{
\hypersetup{linkcolor=}
\setcounter{tocdepth}{3}
\tableofcontents
}
\listoffigures
\listoftables

\printnoidxglossaries

\newpage{}

\subsection*{Disclaimer}\label{disclaimer}

These materials do not constitute a formal publication and are for
information only. They are in a pre-review, pre-decisional state and
should not be formally cited or reproduced. They are to be considered
provisional and do not represent any determination or policy of NOAA or
the Department of Commerce.

\newpage{}

Please cite this publication as:

Cope, J. M., and Wetzel, C. R. 2027. Status and Catch Limit Estimation
of Category 3 stocks off the U.S. West Coast in 2027. Prepared by
{[}COMMITTEE{]}. \pageref*{LastPage}{} pp.

\newpage{}

\section{Executive Summary}\label{sec-exec-sum}

\subsection{Assessment Model}\label{sec-assessment-model}

\subsection{Reference Points, Stock Status, and
Projections}\label{sec-reference-points-stock-status-and-projections}

\newpage{}

\section{Introduction}\label{sec-intro}

There are 11 stocks that remain in the Groundfish Fishery Management
Council that do not have benchmark stock assessments and who catch
limits were established in 2011 using catch-only methods (Dick and
MacCall 2010). These are still considered Category 3 stocks, but this
analysis provides an updated estimate of catch limits that use both
catch and length compositions time series (Rudd et al. 2021). This
approach allows data (length compositions) to inform stock status and is
an improvement from catch-only methods that use stock status as and
input. This work approach is able to provide stock assessment output
similar to what standard benchmark assessments provide, but is only used
to inform catch limits and general stock status. If stock status is
below limit reference points, this approach could also provide a trigger
to move toward an assessment that uses additional data (i.e, Category 2
or 1 stock designations).

\subsection{Stock Description, Ecology, and Fishery
Descriptions}\label{sec-stock-stuff}

\subsubsection{\texorpdfstring{Bank Rockfish (\emph{Sebastes
rufus})}{Bank Rockfish (Sebastes rufus)}}\label{bank-rockfish-sebastes-rufus}

Bank Rockfish is a medium-sized rockfish that ranges from California to
British Columbia, Canada, with a core range from northern California to
Southern Orgeon (Love et al. 2002). They reside in mainly deeper shelf
and shallow slope depths and on a variety of bottom features such as
high relief reefs, rock and boudlers, and cobble and mud, but often
affiliated with hard structures (Love et al. 2002). They can be found
both apart from other Bank Rockfish but also in aggregation just above
the sea floor (Love et al. 2002).

The species name is Latin for ``red'', and Bank Rockfish has often been
refered to as ``red widow'' when landed in fisheries. Bank Rockfish are
caught in trawl and fixed gear commerical fisheries. The peak catches
were experienced in the 1980s. Since 2000, commericial catches have been
very low. Recreational catches are\ldots{}

\subsubsection{\texorpdfstring{Redstripe Rockfish (\emph{Sebastes
proriger})}{Redstripe Rockfish (Sebastes proriger)}}\label{redstripe-rockfish-sebastes-proriger}

Redstripe Rockfish is a medium-sized rockfish and has been recorded from
Mexico to the Aluetian Islands, with a core range from central Oregon to
southeast Alaska (Love et al. 2002). They inhabit mostly deeper shelf
shallower slope depths and are found around high-reflief structures in
solo and in small or larger groups. Schooling behavior is just off the
bottom or above structures (Love et al. 2002).

The species name ``proriger'' in Latin means ``prow-bearing'' and refers
to the prominant sympheseal knob and protruding lower jaw, similar to
Pacific Ocean Perch (\emph{S. alutus}). Commercial catches of Restripe
Rockfish has general been low, but has increased substantially over the
last 10 years (and since the last time catch limits were calculated).

\subsubsection{\texorpdfstring{Rosethorn Rockfish (\emph{Sebastes
helvomaculatus})}{Rosethorn Rockfish (Sebastes helvomaculatus)}}\label{rosethorn-rockfish-sebastes-helvomaculatus}

Rosethorn Rockfish are small- to medium-sized ranging from Mexico to
Alaska, with a core distribution from northern California to southeast
Alaska (Love et al. 2002). They are found in shelf to shallow slope
depths and affiliated with boulders and hard structures on muddy
bottoms. They rarely stray far from the bottom and are usually found
hiding near structures, and do not seem to school (Love et al. 2002).

Their distinctive spots (the species name in Latin means
``tawny-spotted'') make identification easier. They are not a target for
any fishery, but are found in trawl and recreational fishery catch in
small numbers. Some of the largest catches occured in the 1990s, but has
generally dropped since 2000.

\subsubsection{\texorpdfstring{Shortraker Rockfish (\emph{Sebastes
borealis})}{Shortraker Rockfish (Sebastes borealis)}}\label{shortraker-rockfish-sebastes-borealis}

Shortraker Rockfish is a very large and very long-lived rockfish, with
its Latin name, meaning ``northern'', indicating that its range is
largely in the northern latitudes. It has been identified from central
California to the Russia, but is most common from British Colubia to
Kamchatka (Love et al. 2002). Thus the portion of the population being
considered here is in the far southern portion of its range. Shortraker
Rockfish are deeper-dwelling as well, found mostly in slope to deep
slope depths and usually in soft substratum with boulders and stones.

Shortraker are taken in trawl and long-line fisheries where they can be
a target and desired stock in some parts of its range, but off the U.S.
west coast are typically incidentially caught.

\subsubsection{\texorpdfstring{Silvergary Rockfish (\emph{Sebastes
brevispinis})}{Silvergary Rockfish (Sebastes brevispinis)}}\label{silvergary-rockfish-sebastes-brevispinis}

Silvergray Rockfish is a larger rockfish with a reported broad range
from Mexico to the Gulf of Alaska, but with a core range of Oregon to
Alaska (Love et al. 2002). They inhabitat shelf to shallow slope depths,
and aggregate in schools with other rockfishes over rocky bottoms (Love
et al. 2002). They have a similar look as Bocaccio (\emph{S.
pauscispinus}) as well as similar wormy internal disposition. The Latin
name means ``short-spinned'' whereas the Latin name for Bocaccio means
``few spines''.

Silvergray Rockfish are often taken in midwater trawls and hook-and-line
fisheries. They are not currently a target in U.S. west coast fisheries,
and their wormy meat makes them less desireable, but they had
substantial removals from the 1940s through the 1990s. Since 2000 the
catches are significantly decreased, though have shown some growth in
the last 10 years.

\subsubsection{\texorpdfstring{Starry Rockfish (\emph{Sebastes
constellatus})}{Starry Rockfish (Sebastes constellatus)}}\label{starry-rockfish-sebastes-constellatus}

Starry Rockfish (constellatus meaning ``starry'') is a small- to medium-
sized rockfish found from northern California down to southern Baja
Calfornia (Love et al. 2002), but more common south of central
California. They are a nearshore to shallow shelf species and usually
found among rocks, boulders, caves and other hard structures they can
find refuge in (Love et al. 2002). They generally are solitary and stick
near the bottom and close to structure.

Starry rockfish size and colorful look do make them attractive to
premium fish markets as well as recreational fisheries in California
despite being smaller. While historical removals demonstarte starry as a
commonly caught species, catches have greatly decreased since 2000.

\subsubsection{\texorpdfstring{Stripetail Rockfish (\emph{Sebastes
saxicola})}{Stripetail Rockfish (Sebastes saxicola)}}\label{stripetail-rockfish-sebastes-saxicola}

Stripetail Rockfish is a small- to medium-sized rockfish known for its
relatively large-sized eyes. They found from Mexico to southern Alaska
(Love et al. 2002). There core distribution is southern California to
British Columbia and within shelf depths. Despite their Latin name,
which means ``rock-dweller'', Stripetail Rockfish are mostly found in
muddy or sandy bottoms with stones or shells (Love et al. 2002).

Stripetail are mostly taken in commercial trawl fisheries and often
discarded because of their smaller sizes. Total removals have increased
over the past 10 years.

\subsubsection{\texorpdfstring{Yellowmouth Rockfish (\emph{Sebastes
reedi})}{Yellowmouth Rockfish (Sebastes reedi)}}\label{yellowmouth-rockfish-sebastes-reedi}

Theh Yellowmouth Rockfish is a larger rockfish that ranges from northern
California into the Gulf of Alaska, but are most commonly occurring form
southern Oregon to southeast Alaska. They are found in deeper shelf and
shallow slope depths among high relief rocks, where they are often found
schooling just above the hard structures.

Yellowmouth are caught in midwater and bottom trawls, commonly with
Pacific Ocean Perch (\emph{S. alutus}). The Latin name refers to a
former Chairman of the Fisheries Research Board of Canada (Love et al.
2002), and is indicative that Yellowmouth Rockfish are a notable fishery
off of British Columbia.

\subsubsection{\texorpdfstring{Flathead Sole (\emph{Hippoglossoides
elassodon})}{Flathead Sole (Hippoglossoides elassodon)}}\label{flathead-sole-hippoglossoides-elassodon}

Flathead Sole range from central California to Russia, with a core west
distribution that starts in Oregon and goes northward {[}Love2011{]}.
They are generally found in mid-shelf to shallow slope depths and live
on sandy to muddy bottoms.

Flathead Sole do not seem to be major targets in trawl fisheries, with
removals not showing any indetifiable trends. Discards can be as much or
more as the landed amounts.

\subsubsection{\texorpdfstring{Pacific Sanddab (\emph{Citharichthys
sordidus})}{Pacific Sanddab (Citharichthys sordidus)}}\label{pacific-sanddab-citharichthys-sordidus}

\subsubsection{\texorpdfstring{Pacific Cod (\emph{Gadus
macrocephalus})}{Pacific Cod (Gadus macrocephalus)}}\label{pacific-cod-gadus-macrocephalus}

Pacific Cod have one of the largest ranges of the fishes being
considered in this analysis. It also have a large history of being
assessed in Alaska ans elsewhere. While its range is recorded from
California to the Yellow Sea, its core population is found north of
Washington State. This analysis thus deals with a very small part of
both range and overall Pacific Cod population. That said, its removals
are some of the largest of all of the fishes in this analysis as well.
So while it is not a major west coast fishery target nor does this
represent a major portion of the overall species range, Pacific cod are
definitely a presence in west coast groundfish fisheries, not often
discarded, are in the groundfish FMP, and thus need an updated catch
limit estimation.

\subsection{Management History}\label{sec-management-history}

\newpage{}

\section{Data}\label{sec-data}

\subsection{Life History}\label{sec-life-history}

\subsection{Catch}\label{sec-catch}

\subsection{Indices and
Standardization}\label{sec-indices-and-standardization}

\subsection{Composition Data}\label{sec-composition-data}

\subsection{Absolute Abundance}\label{sec-absolute-abundance}

\subsection{Environmental/Ecosystem Indicator
Data}\label{sec-environmental-ecosystem-indicator-data}

\newpage{}

\section{Assessment}\label{sec-assmt-config}

\subsection{Current Modeling
Approach}\label{sec-current-modeling-approach}

\subsection{Configuration of the Base
Model}\label{sec-configuration-of-the-base-model}

\subsection{Bridging}\label{sec-bridging}

\newpage{}

\subsection{Modeling Results}\label{sec-assmt-results}

\subsubsection{Parameter Estimates}\label{sec-parameter-estimates}

\subsubsection{Time Series}\label{sec-time-series}

\subsubsection{Model Fits}\label{sec-model-fits}

\subsubsection{Model Diagnostics}\label{sec-model-diagnostics}

\newpage{}

\subsection{Sensitivity Analyses}\label{sec-assmt-sens}

\newpage{}

\subsection{Management Benchmarks}\label{sec-assmt-bench}

\newpage{}

\subsection{Projections}\label{sec-assmt-proj}

\newpage{}

\section{Discussion}\label{sec-discussion}

\newpage{}

\section{Acknowledgements}\label{sec-acknowledgements}

This document was produced using the R package asar (Schiano et al.
2026), which is free to use and publicly available on
\href{https://github.com/nmfs-ost/asar}{GitHub}.

\newpage{}

\section{References}\label{sec-refs}

\phantomsection\label{refs}
\begin{CSLReferences}{1}{0}
\bibitem[\citeproctext]{ref-dickEstimatesSustainableYield2010}
Dick, E.J., and MacCall, A.D. 2010. Estimates of sustainable yield for
50 data-poor stocks in the {Pacifc} coast ground {Fishery Management
Plan}. NOAA Technical Memorandum NMFS \textbf{NOAA-TM-NMFS-SWFSC-460}.

\bibitem[\citeproctext]{ref-Love2002}
Love, M., Yoklavich, M., and Thorsteinson, L. 2002. {The rockfishes of
the northeast Pacific}. University of California Press, Berkeley, CA,
USA.

\bibitem[\citeproctext]{ref-RuddCatchLengthModels2021}
Rudd, M.B., Cope, J.M., Wetzel, C.R., and Hastie, J. 2021. Catch and
{Length Models} in the {Stock Synthesis Framework}: {Expanded
Application} to {Data-Moderate Stocks}. Frontiers in Marine Science
\textbf{8}. Available from
\url{https://www.frontiersin.org/articles/10.3389/fmars.2021.663554}
{[}accessed 25 August 2022{]}.

\bibitem[\citeproctext]{ref-asar_2026}
Schiano, S., Breitbart, S., and Saul, S. 2026. Asar: Build NOAA stock
assessment report. Available from
\url{https://github.com/nmfs-ost/asar}.

\end{CSLReferences}

\newpage{}

\clearpage

\newpage{}

\section{Figures}\label{sec-figures}

Please refer to the \texttt{stockplotr} package downloaded from
remotes::install\_github(`nmfs-ost/stockplotr') to add premade figures.

\newpage{}

\appendix
\setcounter{figure}{0}
\setcounter{table}{0}

\section{Appendices}\label{sec-appendix}




\end{document}
